%TC:ignore
\section{Reflection Perspective}
\begin{comment}
Instead of having a subsection for each operation, maintenance etc, just make a subsection with something relevant like "Going from Vagrant to Terraform". Ideas: ROI when adding (complexity), Docker Swarm (is it really worth introducing added complexity when our entire application could probably run in one container using sqlite), going from vagrant, to bash script, to terraform - Laurits % Database switching (not dropping requests while switching to a new one, also migration of database between VMs), Vagrant (remote_files were always a manual step only able to be done by some select people), Monitoring, (sucked)

Describe the biggest issues, how you solved them, and which are major lessons learned with regards to:

\item evolution and refactoring,
\item operation, and % 
\item maintenance % 
of your ITU-MiniTwit systems. Link back to respective commit messages, issues, tickets, etc. to illustrate these.

\end{comment}
%TC:endignore
\subsection{Resolving Token Exchange Issues in Docker Swarm Setup}

We initially used a combination of Vagrant and shell scripts to provision and configure our VMs. The Vagrantfile handled the creation and provisioning of VMs, while the shell scripts managed environment configuration. 

Integrating Docker Swarm into this workflow came with significant challenges. Our approach was to initiate a Docker Swarm within the shell script invoked by the Vagrantfile, however this introduced difficulties when retrieving the token generated by \texttt{docker swarm init} to connect workers to the manager. 
After exploring several options, we ultimately moved away from the Vagrant-based setup. Instead, we used shell scripts to provision our infrastructure via the Digital Ocean API, which is better maintained and offers robust configuration options, for example it allows SSH keys to be passed as parameters during VM creation. Towards the end of the course, we further improved our deployment process by adopting Terraform. This provided a more efficient IaC solution.

\subsection{Challenges in Monitoring System Stability with Prometheus and Grafana}
As briefly mentioned in \ref{section:monitoring}, monitoring our system was a near constant struggle throughout this project. Introducing Prometheus and Grafana to our stack led to our application crashing several times and our communication channel getting flooded with alerts from Grafana about Prometheus being down. Simultaneously, we were noticing an increased use of CPU resources. We did not think these two to be related initially and we were investigating other potential reasons for the increasing CPU usage, but when we decided to put off on restarting our Prometheus instance, we noticed that the CPU usage of our droplet stayed consistently low (see \cite{noauthor_minitwit_nodate}).

We spent weeks looking into this issue, both by looking for solutions online and by consulting with other groups, but the solutions found did not work for us. Finally, we decided to (1) not expose the Prometheus port to the public to prevent unauthorised people from querying Prometheus data, (2) limit the CPU usage of the Prometheus container to a maximum of 20\%, and (3) increase the scrape interval from 15 seconds to 60 seconds. These three changes seemed to solve the issue, though we did not realise that by doing (3) on the VM directly, that change would get overwritten at next deployment, leading us to puzzle over a Grafana dashboard that sometimes displayed no data due to queries timeing out. To this day, we have not figured out what caused this issue in the first place, as there were groups with an identical setup that did not experience this. We might have identified the issue earlier if our logging setup had been operational sooner, which would have provided valuable insight.

\subsection{Lessons from an Easter Outage}

The day before the start of the Easter break, our production system transferred from being on a single droplet to a Docker Swarm setup. We came back and noticed that our application had been down for the entirety of the break, including the monitoring system responisble for notifying us. Simulator errors doubling was a stress-factor, which led to us understanding why so many companies schedule production freezes.

%TC:ignore
\subsection{DevOps culture}
% Also reflect and describe what was the "DevOps" style of your work. For example, what did you do differently to previous development projects and how did it work?

%TC:endignore

Our DevOps culture was oriented on implementing the Three Ways principles from The DevOps Handbook \cite[Chapter~1]{kim2021devops}. We implemented \textit{left-to-right flow} using GitHub's issue tracker, kanban board and integration of a CI/CD pipeline. \textit{Right-to-left} flow was achieved using monitoring and logging, assuring that we can swarm the problems as soon as they are detected. Having the responsibility of keeping a website continuously up is a full-time job and was quite stressful to manage alongside other courses.

We practiced the \textit{Third Way} in our team through our ways of working. We made use of pair programming, did mental check-ins before we started working and created a culture of trust.
Issues were self-assigned by team members, allowing people to work on what found find interesting, resulting in a self-organised and motivated team.

